% 05-elevation-bands. figures/bands_swe_era5.pdf at full body height, with the
% numbers the figure cannot state on its face beside it.
%
% The figure carries its own title ("Colorado Rocky Mountains ERA5 1981-2026")
% and its own band labels, so the frametitle names the CUT rather than
% repeating them.
%
% HEIGHT. The plot is 1.19:1, so on a 16:9 body it is height-bound, not
% width-bound — widening its column does nothing and only the height figure
% below moves it. 6.55cm is what clears the frametitle bar and the footline
% with the caption column beside it; 6.75 overran the page by 10pt. \centering
% rather than the center environment, whose \topsep glue was part of that.
%
% Every value in the table is from results/FINDINGS.md, "By elevation band",
% which is generated from the daily checkpoints and never hand-edited. Percent
% figures are of that band's own 46-year mean, not of the domain mean. ERA5
% only: MERRA-2's 1 April value is degenerate in this domain and is excluded
% by the claims discipline (CLAUDE.md), which is also why the aside says which
% model the numbers are from.
%
% It says it PLAINLY. This aside used to read "every number names its product",
% which is not a sentence anyone in the field writes: "data product" is ordinary
% Earthdata vocabulary, but the aphorism that makes the NUMBER the actor was
% invented here and had spread to the README and a skill before it was caught.
% State the source dataset the way a caption would.
\begin{frame}[t]
\frametitle{ERA5 by Elevation Band}
\vspace{0.15em}
\begin{columns}[T,onlytextwidth]
\begin{column}{0.575\textwidth}
{\centering\includegraphics[height=6.55cm]{bands_swe_era5.pdf}\par}
\end{column}
\begin{column}{0.405\textwidth}
\vspace{0.75cm}
\begin{block}{1 April SWE, per cent of band mean}
\scriptsize
\renewcommand{\arraystretch}{1.35}
\begin{tabular}{@{}l@{\hspace{0.9em}}r@{\hspace{0.9em}}r@{}}
                    & \textbf{2023} & \textbf{2026}\\
8,000--14,500\,ft   & 140\%         & 51\%\\
6,500--8,000\,ft    & 297\%         & 5\%\\
\end{tabular}
\end{block}
\vspace{0.40em}
{\raggedright\scriptsize\color{accentDeep}
WY2026 is the lowest 1 April SWE of the 46 years, at 41\% of the record mean.
The deficit is not uniform: the lower band all but vanished.\par}
\vspace{0.55em}
\figaside[\raggedright]{All values ERA5. MERRA-2 gives different amounts for
the same snow.}
\end{column}
\end{columns}
\end{frame}
